% =============================================================================
% equiparticion_y_virial.tex — Nota teórica complementaria: de dónde salen
% los canales alpha_TB y alpha_ΔVB. Equipartición, virial, sólidos armónicos,
% y el origen anarmónico del exceso Delta V_B - T_B. Con historia y fuentes
% primarias/textos de referencia para cada resultado.
% Compilar: pdflatex equiparticion_y_virial.tex (dos pasadas)
% =============================================================================
\documentclass[11pt,a4paper]{article}

\usepackage[utf8]{inputenc}
\usepackage[T1]{fontenc}
\usepackage[spanish,es-nodecimaldot]{babel}
\usepackage[margin=2.3cm]{geometry}
\usepackage{amsmath,amssymb}
\usepackage{booktabs}
\usepackage{enumitem}
\usepackage{xcolor}
\usepackage{graphicx}
\usepackage[colorlinks=true,linkcolor=blue!50!black,citecolor=blue!50!black]{hyperref}

\newcommand{\eps}{\varepsilon}
\newcommand{\kb}{k_B}
\newcommand{\lammps}{\textsc{Lammps}}
\newcommand{\clave}[1]{\textcolor{red!55!black}{\textbf{#1}}}

\title{\textbf{Calor cinético y calor potencial en un cúmulo que vibra}\\
\large Equipartición, el teorema del virial, y por qué
$T_B \approx \Delta V_B$ --- historia, derivaciones y fuentes}
\author{Felipe Wasaff}
\date{Julio 2026}

\begin{document}
\maketitle

\begin{abstract}
Cuando la bola de nuestro proyecto se calienta al chocar, esa energía se
reparte en dos mitades casi iguales: una cinética ($T_B$, los átomos se
mueven) y una potencial ($\Delta V_B$, los enlaces se estiran y comprimen).
Este documento explica \emph{por qué} ocurre eso, con las herramientas que
lo predicen ---el teorema de equipartición y el teorema del virial---, su
historia, sus demostraciones, y por qué la igualdad no es exacta sino que
tiene un exceso sistemático que crece con la energía del impacto. Todas las
afirmaciones históricas y las fórmulas están referenciadas a fuentes
primarias o a textos de posgrado estándar.
\end{abstract}

\tableofcontents

% =============================================================================
\section{La pregunta}
% =============================================================================

En la Sección~\ref{sec:proyecto} definimos dos canales de la partición de
energía del rebote:
\[ \alpha_{T_B} = \frac{T_B}{E_0}, \qquad
   \alpha_{\Delta V_B} = \frac{\Delta V_B}{E_0}, \]
y observamos, en 210 simulaciones independientes, que
$\Delta V_B \approx T_B$ ---con un exceso sistemático,
$\Delta V_B - T_B$, que crece con la velocidad de impacto. Esa igualdad
aproximada no es un accidente numérico: es la predicción de dos teoremas
centenarios de la física estadística y la mecánica clásica, que se aplican
aquí con una precisión notable porque nuestra simulación es
\emph{clásica y sin aproximaciones de temperatura}: no hay efectos
cuánticos que corregir, así que la predicción se cumple de manera casi
exacta, no como un límite de alta temperatura.

% =============================================================================
\section{El teorema de equipartición}
\label{sec:equiparticion_teo}
% =============================================================================

\subsection{Enunciado}

En el ensamble canónico clásico (temperatura $T$ fija, distribución de
Boltzmann $\rho \propto e^{-H/\kb T}$), si una variable $x$ del
Hamiltoniano $H$ tiene un rango de integración tal que
$x\,e^{-H/\kb T}\to0$ en sus extremos (siempre cierto para momentos, y para
coordenadas confinadas por un potencial), entonces
\begin{equation}
\left\langle x\,\frac{\partial H}{\partial x}\right\rangle = \kb T.
\label{eq:equip_general}
\end{equation}

\subsection{Demostración}

Sea $\beta=1/\kb T$. Como
$(\partial H/\partial x)\,e^{-\beta H} = -\beta^{-1}\,\partial_x
\big(e^{-\beta H}\big)$,
\[ \left\langle x\frac{\partial H}{\partial x}\right\rangle
   = -\frac{1}{\beta Z}\int x\,\partial_x\!\big(e^{-\beta H}\big)\,d\Gamma. \]
Integrando por partes en $x$ (el término de borde se anula por hipótesis):
\[ \int x\,\partial_x\!\big(e^{-\beta H}\big)\,dx
   = \big[x\,e^{-\beta H}\big]_{\rm borde} - \int e^{-\beta H}\,dx
   = -\int e^{-\beta H}\,dx, \]
de modo que $\langle x\,\partial H/\partial x\rangle
= \beta^{-1} = \kb T$. $\blacksquare$

\subsection{Corolario: energía por grado de libertad cuadrático}

Si $H$ contiene un término $H\supset a\,x^2$ (con $a$ constante), entonces
$x\,\partial H/\partial x \supset 2a x^2$, y la ecuación~(\ref{eq:equip_general})
da directamente
\begin{equation}
\langle a\,x^2 \rangle = \tfrac12 \kb T.
\label{eq:equip_cuadratico}
\end{equation}
Esta es la forma familiar: \emph{cada grado de libertad que entra al
cuadrado en la energía aporta $\tfrac12\kb T$ en promedio}, sin importar la
masa ni la constante de resorte. Para energía cinética
$K=\tfrac12 m v^2$, cada componente de velocidad aporta $\tfrac12\kb T$;
para los $3N-3$ grados cinéticos internos de nuestra bola,
\begin{equation}
\langle T_B\rangle = \frac{3N-3}{2}\,\kb T
\quad\Longrightarrow\quad
T_{\rm ef} \equiv \frac{2\,T_B}{3N-3}
\label{eq:Tef_derivada}
\end{equation}
---exactamente la definición de temperatura efectiva que usamos en el
proyecto.

\subsection{Historia}

\begin{itemize}[leftmargin=1.4em]
\item \textbf{John James Waterston (1845)} propuso una versión temprana del
reparto de energía en un artículo enviado a la Royal Society. Fue
\clave{rechazado} por los revisores, que lo consideraron especulativo y sin
fundamento suficiente. El manuscrito quedó archivado.
\item \textbf{Maxwell (1860)} llegó al mismo resultado de forma
independiente para la teoría cinética de los gases, repartiendo energía
entre los grados de libertad traslacionales.
\item \textbf{Boltzmann (1871)} lo generalizó a sistemas mecánicos
arbitrarios, más allá de los gases ideales.
\item \textbf{Rayleigh}, revisando los archivos de la Royal Society casi
medio siglo después, encontró el manuscrito rechazado de Waterston y lo
hizo publicar \textbf{47 años después} de haber sido escrito --- un caso
célebre de revisión por pares fallida en la historia de la física.
\item \textbf{Tolman (1918)} dio la demostración general moderna, la que
seguimos arriba, en el lenguaje del espacio de fases.
\item Nota histórica adicional: aplicado ingenuamente a los modos del
campo electromagnético (radiación de cuerpo negro), este mismo teorema
predice energía infinita por unidad de volumen --- la ``catástrofe
ultravioleta''. Ese fracaso fue una de las motivaciones centrales que
llevaron a Planck (1900) y Einstein (1905) a introducir la cuantización.
\end{itemize}

\subsection{Fuentes}
\begin{itemize}[leftmargin=1.4em]
\item J.J. Waterston, \textit{Phil. Trans. R. Soc.} \textbf{183}, 1 (1892)
--- el artículo de 1845, publicado póstumamente gracias a Rayleigh.
\item R.C. Tolman, \textit{Phys. Rev.} \textbf{11}, 261 (1918) --- prueba
general moderna.
\item F. Reif, \textit{Fundamentals of Statistical and Thermal Physics}
(McGraw-Hill, 1965), cap.~6 --- la derivación pedagógica estándar de
posgrado.
\item L.D. Landau y E.M. Lifshitz, \textit{Statistical Physics}, Curso de
Física Teórica vol.~5 (Pergamon) --- tratamiento riguroso alternativo.
\end{itemize}

% =============================================================================
\section{El teorema del virial}
\label{sec:virial_teo}
% =============================================================================

La equipartición predice cuánto vale $T_B$. Pero \emph{no} predice, por sí
sola, que $\Delta V_B$ valga lo mismo --- eso requiere un teorema distinto:
el virial.\footnote{Este es el virial \emph{mecánico} de Clausius,
$2\langle K\rangle=-\langle\sum_i\vec F_i\cdot\vec r_i\rangle$, sobre
trayectorias acotadas en el tiempo. No debe confundirse con la
\emph{expansión del virial} de la mecánica estadística de gases reales,
$pV/Nk_BT=1+B_2(T)N/V+\cdots$: ambas llevan el nombre de Clausius (que
introdujo el término en 1870) pero son resultados distintos, uno mecánico
y el otro estadístico.}

\subsection{Demostración}

Sea $G = \sum_i \vec p_i \cdot \vec r_i$ (suma sobre partículas). Su
derivada temporal es
\[ \frac{dG}{dt} = \sum_i \left(\dot{\vec p}_i\cdot\vec r_i
   + \vec p_i\cdot\dot{\vec r}_i\right)
   = \sum_i \vec F_i\cdot\vec r_i + \sum_i \frac{p_i^2}{m_i}
   = \sum_i \vec F_i\cdot\vec r_i + 2K. \]
Si el movimiento está acotado (posiciones y momentos permanecen finitos,
como en un sólido o una órbita), el promedio temporal de $dG/dt$ sobre un
intervalo largo $\mathcal T$ tiende a cero:
$\langle dG/dt\rangle = [G(\mathcal T)-G(0)]/\mathcal T \to 0$. Por lo
tanto
\begin{equation}
2\langle K\rangle = -\Big\langle\sum_i \vec F_i\cdot\vec r_i\Big\rangle
\equiv \langle \text{virial}\rangle.
\label{eq:virial_general}
\end{equation}
Clausius llamó a $-\sum_i\vec F_i\cdot\vec r_i$ el \emph{virial} del
sistema (del latín \textit{vis}, fuerza).

\subsection{Forma para potenciales tipo ley de potencia}

Si una partícula se mueve en un potencial central $V(r)=c\,r^n$, entonces
$\vec F\cdot\vec r = -nc\,r^n = -nV(r)$, y la
ecuación~(\ref{eq:virial_general}) se reduce a
\begin{equation}
2\langle K\rangle = n\,\langle V\rangle.
\label{eq:virial_potencia}
\end{equation}

\textbf{Ejemplo 1 --- resorte ($n=2$):} $2\langle K\rangle=2\langle V\rangle
\Rightarrow \langle K\rangle=\langle V\rangle$. Esta es la relación que
predice $T_B\approx\Delta V_B$ en nuestro sistema, una vez que aproximamos
el potencial de Lennard-Jones cerca de su mínimo como una parábola
(Sección~\ref{sec:solido_armonico}).

\textbf{Ejemplo 2 --- gravedad o Coulomb ($n=-1$):}
$2\langle K\rangle=-\langle V\rangle$. Esta forma del virial tiene una
aplicación histórica célebre: \textbf{Zwicky (1933)}, midiendo las
velocidades de las galaxias en el cúmulo de Coma y aplicando
$2\langle K\rangle=-\langle V\rangle$ para estimar la masa total del
cúmulo, encontró que la masa necesaria para mantenerlas gravitacionalmente
ligadas superaba largamente la masa visible en estrellas --- el primer
indicio observacional de lo que hoy llamamos materia oscura.

\subsection{Historia}

\textbf{Rudolf Clausius (1870)} introdujo el teorema y el término
``virial'' en el contexto de \emph{corregir la ley de los gases ideales}:
quería entender cómo las fuerzas entre moléculas reales modifican la
relación presión-volumen-temperatura. Ese origen sigue vivo hoy en el
nombre de las \emph{expansiones viriales} de la ecuación de estado,
$pV = nRT(1 + B_2/V + B_3/V^2 + \cdots)$, usadas en termodinámica de gases
reales.

\subsection{Fuentes}
\begin{itemize}[leftmargin=1.4em]
\item R. Clausius, \textit{Ueber einen auf die Wärme anwendbaren
mechanischen Satz}, \textit{Ann. Phys.} \textbf{141}, 124 (1870) ---
artículo original.
\item F. Zwicky, \textit{Die Rotverschiebung von extragalaktischen Nebeln},
\textit{Helv. Phys. Acta} \textbf{6}, 110 (1933) --- aplicación a Coma y
el origen de la ``materia oscura''.
\item H. Goldstein, C. Poole y J. Safko, \textit{Classical Mechanics},
3.\ ed. (Addison-Wesley, 2002), cap.~3 --- derivación clásica estándar de
posgrado (la que seguimos arriba).
\end{itemize}

% =============================================================================
\section{Un sólido que vibra: equipartición + virial juntos}
\label{sec:solido_armonico}
% =============================================================================

Cerca de su mínimo $r_m$, el potencial de Lennard-Jones se puede aproximar
como una parábola (expansión de Taylor a segundo orden):
\begin{equation}
V_{\rm LJ}(r) \approx V(r_m) + \tfrac12 V''(r_m)\,(r-r_m)^2,
\qquad V''(r_m) \approx 57.15\,\eps/\sigma^2
\label{eq:taylor2}
\end{equation}
(el valor numérico es el que calculamos para nuestro sistema en la guía
principal del proyecto). Esta es la aproximación \emph{armónica}: cada par
de átomos vecinos se comporta como un resorte.

Un sólido de $N$ átomos vibrando armónicamente equivale, tras
diagonalizar, a $3N-3$ osciladores independientes (modos normales, en el
lenguaje moderno de fonones). Cada modo tiene, por
equipartición~(\ref{eq:equip_cuadratico}), $\tfrac12\kb T$ de energía
cinética y, por el virial con $n=2$, la misma cantidad de energía
potencial:
\begin{equation}
\langle T_B\rangle = \frac{3N-3}{2}\kb T, \qquad
\langle \Delta V_B\rangle = \frac{3N-3}{2}\kb T
\quad\Longrightarrow\quad
\langle T_B\rangle = \langle \Delta V_B\rangle.
\label{eq:virial_solido}
\end{equation}
Esta es la predicción teórica exacta (a nivel armónico) detrás de lo que
medimos en las 210 simulaciones.

\subsection{Conexión con la ley de Dulong-Petit}

Sumando las dos mitades, cada átomo del sólido almacena $3\kb T$ de
energía vibracional total (3 modos cinéticos + 3 potenciales, cada uno
$\tfrac12\kb T$). Para un mol de átomos, esto da una capacidad calorífica
$C_V = 3R$ --- la ley empírica de \textbf{Dulong y Petit (1819)}, que
describieron un valor casi universal para sólidos a temperatura ambiente
mucho antes de que existiera una explicación microscópica. La
equipartición, medio siglo después, fue la explicación teórica.

\subsection{Por qué no hace falta corrección cuántica aquí}

Dulong-Petit falla en materiales reales a baja temperatura: la capacidad
calorífica cae a cero cuando $T\to0$, en vez de mantenerse constante en
$3R$. La razón es cuántica: los modos de vibración de alta frecuencia
requieren $\kb T \gtrsim \hbar\omega$ para excitarse térmicamente, y a
baja $T$ quedan ``congelados''. \textbf{Einstein (1907)} dio la primera
explicación cuántica (osciladores idénticos de una sola frecuencia);
\textbf{Debye (1912)} la refinó con un espectro continuo de frecuencias,
reproduciendo la ley $T^3$ observada experimentalmente a muy baja
temperatura.

\clave{Nuestra simulación no necesita ninguna de estas dos correcciones.}
\lammps{} integra dinámica clásica de Newton --- no hay $\hbar$, ningún
modo está ``congelado''. El virial clásico se cumple exactamente, no como
límite de alta temperatura sino como la física completa y correcta del
modelo. Por eso la igualdad $T_B\approx\Delta V_B$ es tan limpia en
nuestros datos.

\subsection{Fuentes}
\begin{itemize}[leftmargin=1.4em]
\item P.L. Dulong y A.T. Petit, \textit{Ann. Chim. Phys.} \textbf{10}, 395
(1819) --- ley empírica original.
\item A. Einstein, \textit{Ann. Phys.} \textbf{22}, 180 (1907) --- primer
modelo cuántico de sólido.
\item P. Debye, \textit{Ann. Phys.} \textbf{39}, 789 (1912) --- refinamiento
con espectro de frecuencias.
\item N.W. Ashcroft y N.D. Mermin, \textit{Solid State Physics} (Saunders
College, 1976) --- tratamiento estándar de posgrado de todo lo anterior.
\end{itemize}

% =============================================================================
\section{El exceso anarmónico: por qué $\Delta V_B$ es \emph{un poco más}
que $T_B$}
\label{sec:anarmonia}
% =============================================================================

La ecuación~(\ref{eq:taylor2}) es solo el segundo orden de una expansión
de Taylor. El siguiente término, cúbico,
\begin{equation}
V_{\rm LJ}(r) \approx V(r_m) + \tfrac12 V''(r_m)(r-r_m)^2
+ \tfrac16 V'''(r_m)(r-r_m)^3 + \cdots,
\end{equation}
\textbf{rompe la simetría} de la parábola. En un potencial puramente
armónico (solo el término cuadrático), la probabilidad de encontrar al
átomo desplazado $+u$ o $-u$ de su posición de equilibrio es idéntica
(simetría par), así que el desplazamiento promedio es exactamente cero:
$\langle u\rangle=0$. Con el término cúbico presente, esa simetría se
rompe: el pozo de potencial es más ``duro'' hacia adentro (repulsión
$\propto r^{-12}$, muy empinada) y más ``blando'' hacia afuera (atracción
$\propto r^{-6}$, mucho más suave). El átomo, vibrando con cierta energía,
pasa más tiempo y alcanza mayor amplitud hacia el lado barato ---el
promedio $\langle u\rangle$ deja de ser cero y se vuelve positivo: la
distancia promedio entre átomos \textbf{crece con la energía vibracional}.

Esto \textbf{es}, microscópicamente, la dilatación térmica. La conexión
sistemática entre la forma (asimetría) del potencial interatómico y la
dilatación térmica fue establecida por \textbf{Grüneisen (1912)}, quien
introdujo el parámetro que lleva su nombre,
$\gamma = -\dfrac{d\ln\omega}{d\ln V}$, cuantificando cuánto se ablandan
los modos de vibración cuando el sólido se expande.

En nuestros datos, este exceso se ve directamente en el cociente
$\alpha_{\Delta V_B}/\alpha_{T_B}$, que crece de $\approx1.1$
($v_{\rm imp}=1.16\,\sigma/\tau$) a $\approx1.7$
($v_{\rm imp}=2.02\,\sigma/\tau$): a mayor energía de impacto, mayor
amplitud de vibración, y mayor el sesgo hacia el lado ``caro'' de la
asimetría de Lennard-Jones.

\subsection{Fuentes}
\begin{itemize}[leftmargin=1.4em]
\item E. Grüneisen, \textit{Theorie des festen Zustandes einatomiger
Elemente}, \textit{Ann. Phys.} \textbf{344}, 257 (1912) --- origen del
parámetro de Grüneisen y la conexión anarmonicidad--dilatación térmica.
\item N.W. Ashcroft y N.D. Mermin, \textit{Solid State Physics}, cap.\ de
efectos anarmónicos --- tratamiento moderno estándar.
\end{itemize}

% =============================================================================
\section{Aplicación explícita a este proyecto}
\label{sec:proyecto}
% =============================================================================

Reuniendo todo, las cantidades que medimos en la simulación y sus
predicciones teóricas:

\begin{center}
\begin{tabular}{lll}
\toprule
Cantidad & Definición medida & Predicción teórica \\
\midrule
$T_{\rm ef}$ & $2T_B/(3N-3)$ & equipartición, ec.~(\ref{eq:Tef_derivada}) \\
$T_B \approx \Delta V_B$ & tabla partición, Sec.~\ref{sec:solido_armonico} & virial armónico, ec.~(\ref{eq:virial_solido}) \\
$\Delta V_B - T_B > 0$, crece con $v$ & Sec.~\ref{sec:anarmonia} & anarmonicidad (Grüneisen) \\
\bottomrule
\end{tabular}
\end{center}

Las dos secciones siguientes contrastan estas predicciones contra el sweep
con orientaciones (\texttt{data/outputs/terminado/cluster\_orient}): $N\in\{1,2,7,13,
55,79,87\}\times$ 6 alturas $\times$ 5 orientaciones iniciales del cúmulo
$=210$ simulaciones independientes, las mismas que se citan en el
resumen. Cada punto reportado abajo es media $\pm$ desviación estándar
sobre las 5 orientaciones a igual $(N,h_0)$ --- no un único run, a
diferencia de las cifras que aparecían en una versión anterior de esta
nota.

\subsection{¿Se cumple $T_B\approx\Delta V_B$? Contraste directo}
\label{sec:contraste_directo}

La Fig.~\ref{fig:virial_check}a muestra $\alpha_{T_B}$ y
$\alpha_{\Delta V_B}$ en función de $v_{\rm imp}$ para $N=79$, el cúmulo
canónico del proyecto:

\begin{center}
\begin{tabular}{cccc}
\toprule
$v_{\rm imp}\,(\sigma/\tau)$ & $\alpha_{T_B}$ & $\alpha_{\Delta V_B}$ & $\alpha_{\Delta V_B}/\alpha_{T_B}$\\
\midrule
0.70 & $0.13\pm0.09$ & $0.04\pm0.02$ & $0.98\pm1.08$ \\
0.93 & $0.17\pm0.05$ & $0.21\pm0.17$ & $1.12\pm0.66$ \\
1.16 & $0.22\pm0.10$ & $0.24\pm0.10$ & $1.10\pm0.17$ \\
1.43 & $0.21\pm0.06$ & $0.26\pm0.09$ & $1.24\pm0.31$ \\
1.72 & $0.19\pm0.04$ & $0.31\pm0.05$ & $1.65\pm0.36$ \\
2.02 & $0.18\pm0.02$ & $0.31\pm0.05$ & $1.68\pm0.22$ \\
\bottomrule
\end{tabular}
\end{center}

Las dos curvas de la Fig.~\ref{fig:virial_check}a son estadísticamente
indistinguibles ---dentro de una barra de error--- en la ventana
$0.9\lesssim v_{\rm imp}\lesssim1.4$: ahí el virial armónico,
ec.~(\ref{eq:virial_solido}), \textbf{se cumple}. Fuera de esa ventana pasan
dos cosas distintas, cada una con una explicación física propia (no son
``ruido''): a $v$ alta $\alpha_{\Delta V_B}$ se despega sistemáticamente
por encima de $\alpha_{T_B}$ (Sec.~\ref{sec:exceso_v_alta}); a $v$ baja
($v_{\rm imp}=0.70$) ambas fracciones son pequeñas y sus barras de error
son comparables a su magnitud, señal de que el régimen de contacto es
cualitativamente distinto (Sec.~\ref{sec:falla_v_baja}).

\begin{figure}[h!]
\centering
\includegraphics[width=\textwidth]{../figuras_rbef/fig_virial_check.pdf}
\caption{Contraste del virial armónico contra el sweep con orientaciones.
(a) $\alpha_{T_B}$ y $\alpha_{\Delta V_B}$ vs.\ $v_{\rm imp}$ para $N=79$
(media $\pm\sigma$, 5 orientaciones por punto): las curvas coinciden dentro
de error para $0.9\lesssim v_{\rm imp}\lesssim1.4$ y se separan
sistemáticamente a $v$ mayor. (b) Razón $\alpha_{\Delta V_B}/\alpha_{T_B}$
para $N=13,55,79,87$: el mismo patrón ---acercamiento a la línea virial
(razón $=1$) en una ventana intermedia de $v$, seguido de un exceso
creciente--- se repite de forma robusta en los cuatro tamaños de cúmulo
``sólidos'' del sweep.}
\label{fig:virial_check}
\end{figure}

\subsection{El exceso a $v$ alta: anarmonicidad, y que no depende de $N$}
\label{sec:exceso_v_alta}

La Fig.~\ref{fig:virial_check}b generaliza la razón
$\alpha_{\Delta V_B}/\alpha_{T_B}$ a $N=13,55,79,87$. El resultado clave es
que el mismo patrón cualitativo aparece en los cuatro tamaños: la razón
cruza $1$ alrededor de $v_{\rm imp}\approx1.1$--$1.4$ y luego crece de
forma monótona, alcanzando $1.6$--$1.8$ en el extremo del sweep
($v_{\rm imp}\approx2.0$). Que la tendencia sea la misma en un rango de
$N$ que va de 13 a 87 átomos --- con geometrías, números de coordinación y
frecuencias de modo normal distintas entre sí --- es evidencia de que el
exceso no es un artefacto de tamaño finito ni de la definición por cierre
de $\alpha_{\Delta V_B}$ (Sec.~\ref{sec:anarmonia}), sino la firma genuina
de la anarmonicidad de Lennard-Jones predicha por Grüneisen: a mayor
energía de impacto, mayor amplitud de vibración, y el pozo asimétrico
empuja $\langle\Delta V_B\rangle$ por encima de $\langle T_B\rangle$ de
forma sistemática y reproducible entre cúmulos de tamaño muy distinto.

\subsection{Por qué el virial deja de cumplirse a $v$ baja}
\label{sec:falla_v_baja}

La demostración de la Sec.~\ref{sec:virial_teo} depende de una hipótesis
que es fácil pasar por alto: el promedio temporal de $dG/dt$ tiende a cero
solo si el intervalo de promediado $\mathcal T$ es largo comparado con el
período característico del movimiento acotado. En nuestro sistema,
$\mathcal T$ es el \emph{tiempo de contacto} bola-superficie durante el
primer rebote --- no un intervalo elegido por nosotros, sino impuesto por
la dinámica.

A $v_{\rm imp}=0.70$ ($N=79$), $e_1\approx0.92$ (Sec.~\ref{sec:proyecto}):
el impacto es casi elástico, la bola apenas se deforma, y el tiempo de
contacto es corto en comparación con el período de vibración interna de
los enlaces LJ ($\sim2\pi\sqrt{m/V''(r_m)}\approx0.83\,\tau$, a partir de
la ec.~(\ref{eq:taylor2})). El promedio temporal que exige el virial
simplemente no alcanza a completarse dentro del contacto: la bola no
``olvida'' su fase inicial de vibración antes de despegar, y por eso
$\alpha_{T_B}$ y $\alpha_{\Delta V_B}$ son pequeños, ruidosos entre
orientaciones (la fase inicial de cada orientación importa) y su razón no
converge a ningún valor estable (nótese la barra de error de
$\pm1.08$ sobre un valor medio de $0.98$ en la tabla de
Sec.~\ref{sec:contraste_directo}). Esto no contradice el virial: confirma
que su hipótesis de fondo ---separación de escalas temporales entre el
contacto y la vibración interna--- es la que efectivamente controla cuándo
la predicción se cumple.

\subsection{Conclusión: la teoría se cumple, con dos matices esperables}

En resumen, frente a los datos:
\begin{itemize}[leftmargin=1.4em]
\item La \textbf{equipartición} se cumple por construcción (es la
definición de $T_{\rm ef}$) y no es falsable con estos datos por sí sola;
lo que sí es falsable, y se confirma, es que \emph{la misma} $T_{\rm ef}$
gobierna tanto el canal cinético como el potencial (el virial).
\item El \textbf{virial armónico} ($T_B\approx\Delta V_B$) se cumple
cuantitativamente, dentro de barras de error de $\sim$10--30\%, en la
ventana $0.9\lesssim v_{\rm imp}\lesssim1.4$, y este resultado es robusto
para $N=13,55,79,87$: cuatro cúmulos de tamaño muy distinto obedecen la
misma ley.
\item Fuera de esa ventana, las dos desviaciones observadas
---exceso anarmónico creciente a $v$ alta, ruido y razón inestable a $v$
baja--- tienen explicaciones físicas identificables (Grüneisen; separación
de escalas temporales) y no son fallas del formalismo, sino límites
esperables de su rango de validez.
\end{itemize}
Valores de referencia ($N=79$, plateau $0.9<v\lesssim1.4$):
$\alpha_{T_B}\approx0.17$--$0.22$, $\alpha_{\Delta V_B}\approx0.21$--$0.24$;
el acuerdo entre ambos, cuantitativo dentro de las barras de error, es la
evidencia directa de que la bola termaliza como un sólido casi armónico
durante el contacto, tal como exige la equipartición para poder hablar de
``temperatura'' en un sistema de apenas 79 átomos fuera de equilibrio
global.

% =============================================================================
\section*{Bibliografía completa}
% =============================================================================

\begin{thebibliography}{99}

\bibitem{Waterston1892}
J.J. Waterston, \textit{On the Physics of Media that are Composed of Free
and Perfectly Elastic Molecules in a State of Motion}, Phil. Trans. R.
Soc. \textbf{183}, 1 (1892) [escrito en 1845, rechazado entonces,
publicado póstumamente por gestión de Rayleigh].

\bibitem{Tolman1918b}
R.C. Tolman, \textit{A General Theory of Energy Partition with
Applications to Quantum Theory}, Phys. Rev. \textbf{11}, 261 (1918).

\bibitem{Reif1965b}
F. Reif, \textit{Fundamentals of Statistical and Thermal Physics}
(McGraw-Hill, Nueva York, 1965).

\bibitem{LandauLifshitz}
L.D. Landau y E.M. Lifshitz, \textit{Statistical Physics}, Course of
Theoretical Physics Vol.~5 (Pergamon Press, Oxford, 1980), 3.\ ed.

\bibitem{Clausius1870}
R. Clausius, \textit{Ueber einen auf die Wärme anwendbaren mechanischen
Satz}, Ann. Phys. \textbf{141}, 124 (1870).

\bibitem{Zwicky1933}
F. Zwicky, \textit{Die Rotverschiebung von extragalaktischen Nebeln},
Helv. Phys. Acta \textbf{6}, 110 (1933).

\bibitem{Goldstein}
H. Goldstein, C. Poole y J. Safko, \textit{Classical Mechanics}, 3.\ ed.
(Addison-Wesley, San Francisco, 2002).

\bibitem{DulongPetit1819}
P.L. Dulong y A.T. Petit, \textit{Recherches sur quelques points
importants de la théorie de la chaleur}, Ann. Chim. Phys. \textbf{10}, 395
(1819).

\bibitem{Einstein1907}
A. Einstein, \textit{Die Plancksche Theorie der Strahlung und die Theorie
der spezifischen Wärme}, Ann. Phys. \textbf{22}, 180 (1907).

\bibitem{Debye1912}
P. Debye, \textit{Zur Theorie der spezifischen Wärmen}, Ann. Phys.
\textbf{39}, 789 (1912).

\bibitem{AshcroftMermin}
N.W. Ashcroft y N.D. Mermin, \textit{Solid State Physics} (Saunders
College Publishing, Filadelfia, 1976).

\bibitem{Gruneisen1912}
E. Grüneisen, \textit{Theorie des festen Zustandes einatomiger Elemente},
Ann. Phys. \textbf{344}, 257 (1912).

\end{thebibliography}

\end{document}
